% TOP_PROBLEM: {"id": 30006694, "title": "The Exponential Time Hypothesis for 3-SAT", "category_id": 15, "category": {"id": 15, "name": "computer_science", "display_name": "Computer Science", "description": "Computational complexity, algorithms, and theoretical CS.", "slug": "computer-science", "order_index": 15, "created_at": "2026-07-31T15:26:25.671Z"}, "status": "open"}
% FIRST_POSTED: 2026-09-25
% ENTRY_KIND: statement_only
% !TeX program = lualatex
% Definition and sources only; this file is not an LLM solution attempt.
\documentclass[11pt]{article}
\usepackage[margin=25mm]{geometry}
\usepackage{fontspec}
\setmainfont{Latin Modern Roman}
\usepackage{amsmath,amssymb,mathtools}
\usepackage{unicode-math}
\setmathfont{Latin Modern Math}
\usepackage{xurl,hyperref}
\hypersetup{colorlinks=true,urlcolor=blue}
\setcounter{secnumdepth}{0}
\setlength{\parindent}{0pt}
\setlength{\parskip}{5pt}
\setlength{\emergencystretch}{3em}
\begin{document}
\section*{\texorpdfstring{The Exponential Time Hypothesis for 3-SAT}{The Exponential Time Hypothesis for 3-SAT}}\label{30006694}
\subsection{Definitions and mathematical statement}
A $3$-CNF formula is a conjunction of clauses, each a disjunction of at most three literals in Boolean variables. Let $n$ be the number of variables and $|F|$ the input length. ETH asserts a strictly positive exponential-time rate for deciding satisfiability, conventionally expressed as
\[
 \exists c>0:\quad 3\text{-SAT has no deterministic }O(2^{cn}\operatorname{poly}(|F|))\text{-time algorithm}.
\]
Polynomial input-reading factors and the precise time-bound convention are implicit in the catalogue's abbreviated $2^{cn}$ wording. The parameter is the number of variables, not the numerical size of an encoding or the clause count alone.
\par\smallskip\textit{Formulation and concept references:} \href{https://doi.org/10.1006/jcss.2000.1727}{[S1]}.

\subsection{Short English statement}
Does some positive exponential dependence on the number of variables remain unavoidable for deterministic $3$-SAT algorithms?
\subsection{Sources}
\begin{itemize}
\item[S1]R. Impagliazzo, R. Paturi, J. Comput. System Sci. 62(2):367-375, 2001.
\url{https://doi.org/10.1006/jcss.2000.1727}.
\textit{Formulation source.}
\item[S2]Complexity Theory: Stronger forms of P != NP, Roei Tell.
\url{https://www.cs.toronto.edu/~roei/2026%20complexity/6%20stronger%20lbs.pdf}.
\textit{Further reference; not a proof endorsement.}
\end{itemize}
\end{document}
